\FloatBarrier
\section{Inferring and Testing an Effective Model}
\label{app:model-inference}
\label{sec:combining-measurements}

Sections~\ref{sec:electron-number-control} and~\ref{sec:exchange-virtual-tunneling}
use addition energies and an exchange coupling to describe the device.
This appendix develops a supplementary check: whether a common set of
model parameters accounts for the measurements of those quantities.

Charge-transition positions, excitation gaps and charge-distribution
curves constrain different terms of an effective Hamiltonian.
We now combine such data to estimate its parameters. A separate
measurement can then test a prediction that was not used in the fit.
The example below fits charge data and predicts an exchange splitting;
the spin experiment that tests that prediction is developed in
Section~\ref{sec:measurement-two-spin-exchange}.

\subsection{Relating model parameters to measurements}
\label{sec:measurement-model-predictions}

First choose the retained states and interaction terms, as in
Appendix~\ref{sec:device-hubbard-reduction}. Collect the unknown parameters in
$\boldsymbol\theta$. Each measurement then has a prediction
$f_k(\boldsymbol\theta)$, obtained from the corresponding
charge-boundary or spectroscopy calculation.
Refer all data to matched electrode settings, or include the
measured voltage dependence in the model.

Consider the two-dot Hubbard model of
Appendix~\ref{sec:device-hubbard-reduction}, with one orbital per dot,
equal on-site repulsions $U_1=U_2=U$, and zero detuning. Write
$t_c=|t_{12}|$ and $D=U-V_{12}>0$, and omit Zeeman and direct-exchange
terms. For this example, $\boldsymbol\theta=(t_c,D)$.
The charge-polarization measurement in
Appendix~\ref{app:tunnel-coupling} constrains $t_c$.
Suppose spectroscopy also identifies the two mixed singlet
eigenstates of the matrix in Eq.~\eqref{eq:device-exact-hubbard-check}.
Their energy separation is
\begin{equation}
 \Delta_{\mathrm{ch}}=\sqrt{D^2+16t_c^2}.
 \label{eq:measurement-charge-model-prediction}
\end{equation}
Use these two charge measurements to determine $t_c$ and $D$,
and reserve the spin-oscillation measurement of $J$ for testing
the result. This comparison assumes that the same retained-orbital
model describes the charge sectors used in the measurements.

\subsection{Joint estimation, uniqueness, and uncertainty}
\label{sec:measurement-model-estimation}

Let $y_k$ be the $k$th measured value and $\sigma_k$ the standard
deviation of its measurement error. For independent Gaussian errors, estimate
$\boldsymbol\theta$ by minimizing
\begin{equation}
 \chi^2(\boldsymbol\theta)
   =\sum_k\frac{[y_k-f_k(\boldsymbol\theta)]^2}{\sigma_k^2}.
 \label{eq:measurement-joint-fit}
\end{equation}
In the example, $y_1$ is the measured tunnel coupling and $y_2$
the identified singlet excitation gap; the predictions are
$f_1=t_c$ and $f_2=\sqrt{D^2+16t_c^2}$. In the ideal limit of
negligible uncertainty, these give
$t_c=y_1$ and $D=\sqrt{y_2^2-16y_1^2}$.

Check which parameters the chosen data actually determine. If
only $J$ were measured, its weak-tunneling expression
$J\simeq4t_c^2/D$ would constrain a ratio, leaving many pairs
$(t_c,D)$ possible. The two charge measurements resolve that
ambiguity, but still determine only $U-V_{12}$, not $U$ and
$V_{12}$ separately. To separate those parameters, add charge
measurements whose predicted energies depend on them differently.
If the available data cannot distinguish parameter choices within
their uncertainty, report the constrained combinations.

Estimate the parameter uncertainty by refitting simulated
repetitions generated from the fitted predictions and the
measurement-noise model. In each repetition, use the same sampled
calibration error for all measurements sharing that calibration;
it shifts several inferred energies together. Such correlated errors require
their covariance in the fit, or an explicit fit of the calibration
parameters together with the device parameters. Keep the resulting
joint parameter uncertainty for the prediction test below.

\subsection{Testing predictions on independent data}
\label{sec:uncertainty-model-validation}

Use the fitted parameters to predict a quantity excluded from the
fit, without readjusting them to that measurement. For the two-dot
example, Eq.~\eqref{eq:device-exact-hubbard-check} predicts
\begin{equation}
 J_{\mathrm{pred}}
   =\frac{\sqrt{D^2+16t_c^2}-D}{2}.
 \label{eq:measurement-independent-exchange-prediction}
\end{equation}
Compare this with $J$ obtained from the two-spin exchange measurement
of Section~\ref{sec:measurement-two-spin-exchange}.
Using the exact Hubbard expression avoids attributing an error in
$J\simeq4t_c^2/D$ to the Hamiltonian itself. Recompute
$J_{\mathrm{pred}}$ for the joint parameter estimates obtained
above to determine its uncertainty, including uncertainty in the
operating point. If the measured $J$ had been used to determine
$D$, that same $J$ would no longer be an independent test.

Compare the predicted range with the measured $J$ and its
uncertainty. If they disagree, first check the voltage calibration,
thermal broadening in the charge fit, and repeated scans for drift.
If these checks do not account for a reproducible discrepancy at
matched settings, revisit the Hamiltonian assumptions.
For example, low-lying orbital or valley states seen
in spectroscopy may need to be retained, or direct exchange omitted
from the simple model may contribute to $J$. Include such terms
when supported by the data and test the revised model on another
independent measurement or voltage setting.
